%%
%% Chapter 8: Limitations and Future Work
%%
\chapter{Limitations and Future Work}
\label{ch:08}

This study has several limitations that should be considered when interpreting
the results. We state them plainly --- including the quantitative targets the
work set for itself and did not meet --- and then map the concrete follow-up
tracks that the thesis infrastructure leaves ready to run.

\section{Performance Gaps and Accepted Target Misses}
\label{sec:lim-gaps}

\textbf{Accepted target misses.} The project's pre-registered research plan
(PRD) set explicit performance gates before the main experiments were run, and
the final results miss them. On UCF-Crime the minimum gate was 83\% frame-level
AUC (target range 85--87\%); the achieved headline is $82.5 \pm 0.4$\%, i.e.\
0.5 points below the minimum gate and well below the target range. On
XD-Violence the minimum gate was 80\% AP (target range 82--85\%); the achieved
headline is $78.7 \pm 0.9$\%, i.e.\ 1.3 points below the minimum. We report
these as accepted misses rather than reframing the targets after the fact. The
PRD also set a baseline-reproduction gate --- reproducing RTFM to within one
point of its published 84.30\% UCF-Crime AUC --- as a check on the evaluation
pipeline; the RTFM reproduction study, carried out on the XD-Violence I3D
configuration, reached 65.70\% AP against the published 77.81\% and was
likewise accepted as a miss after a diagnostic investigation ruled out an
evaluation bug and attributed the gap to feature-dimensionality and
training-regime differences (Appendix~\ref{app:a}).

\textbf{Why the gap exists.} Our best UCF-Crime result falls below the
published performance of methods that use end-to-end fine-tuned backbones,
learnable text branches, or temporal modeling specifically designed for
anomaly detection. This gap reflects our architectural choice to freeze all
backbone encoders, which prioritizes computational efficiency and modularity
over peak performance. Unfreezing selected backbone layers or incorporating
temporal transformers could improve absolute numbers but would increase
training cost substantially on our single-GPU setup.

\textbf{Positioning within the published landscape.} Chapter~\ref{ch:07}
(Tables~\ref{tab:sota-full} and~\ref{tab:sota-fair}) places these numbers in
the full published landscape; we summarize the read here rather than repeating
the tables. Within the comparable regime --- frozen features, no text branch,
a single-GPU-class head --- our 78.7\% XD-Violence AP is on par with
RTFM~\cite{tian2021rtfm} (77.81\%) and the I3D variant of
MGFN~\cite{chen2023mgfn} (79.19\%),
and above the founding Sultani et al.~\cite{sultani2018ucfcrime} baseline
(75.41\%). On UCF-Crime our 82.5\% clears Sultani but trails the other
frozen-feature methods (RTFM 84.30\%,
Light-WVAD~\cite{wang2024lightwvad} 84.7\%, UR-DMU~\cite{zhou2023urdmu}
86.97\%, MGFN 86.98\%,
CLIP-TSA~\cite{joo2023cliptsa} 87.58\%),
so our competitiveness within this regime is clearest on XD-Violence AP and
weakest on UCF-Crime AUC. We attribute the residual UCF gap to CLIP-TSA ---
the closest frozen-CLIP analogue --- to the temporal self-attention and
multi-crop aggregation it adds and our minimal head omits, but we do not
isolate this experimentally and leave a controlled ablation to future work. A
heavier budget is also not automatically decisive: the training-free methods
EventVAD~\cite{shao2025eventvad} and LAVAD~\cite{zanella2024lavad} run a 7B
and 13B multimodal LLM yet reach
82.03\% and 80.28\% UCF AUC --- at or below our 82.5\% --- and fall far below
us on XD-Violence AP
(64.04\% / 62.01\% vs 78.7\%), at a far higher inference budget.

\section{Seed Sensitivity on XD-Violence}
\label{sec:lim-seeds}

XD-Violence gated fusion shows elevated seed variance only for the
highest-capacity backbone (SigLIP2 Giant, 2.8\% AP standard deviation); the
other three backbones are stable (near 0.9\%), comparable to UCF-Crime
(maximum 0.4\%). As noted in Chapter~\ref{ch:04}, this variance is sensitive
to the temporal smoothness regularizer, which we disable on XD-Violence; the
residual Giant sensitivity suggests that AP's calibration dependence and
XD-Violence's imbalanced categories still create a challenging optimization
landscape for the largest backbone. Practically, this means single-seed
XD-Violence results for high-capacity backbones should be treated with
caution --- a lesson this thesis internalized by reporting all headline and
ablation numbers as three-seed means --- and future work should investigate
ensemble strategies or variance-reduction techniques for more stable training.

\section{Test-Time Adaptation Scope}
\label{sec:lim-tta}

Our test-time method targets the regime where one modality stays robust; its
gains are concentrated on corruptions that collapse the visual-language stream
(e.g., Gaussian noise) and are near-zero, by construction, where both streams
degrade (motion blur, JPEG compression). The gains concern corrupted-input AUC
only; clean-benchmark AUC is untouched by design, and we make no claim
otherwise.

Two scope caps deserve emphasis. First, the method relies on
\emph{transductive} whole-condition statistics: the routed weight for a
condition is estimated from the entire corrupted test set of that condition
before scoring it. This is a standard evaluation protocol in the TTA
literature, but it is not an online deployment recipe --- a real system sees
its test distribution incrementally, not in advance. Until a streaming variant
exists (Section~\ref{sec:lim-future}), the honest claim class is transductive
robustness recovery, and we make no online-adaptation claim. Second, the
discriminative-spread reliability proxy can over-route on backbones whose
visual stream compresses scores without losing rank (SigLIP2-Base at
high-severity Gaussian noise on one seed) --- the proxy conflates score
compression with rank loss, and a label-free signal that distinguishes the two
is an open problem. Both caps are disclosed wherever the corresponding gains
are reported (Chapter~\ref{ch:06}).

\section{Dataset Scope and Synthetic Corruptions}
\label{sec:lim-scope}

Our evaluation is limited to two datasets (UCF-Crime and XD-Violence).
Generalization to other VAD benchmarks such as ShanghaiTech or CUHK Avenue,
which feature different anomaly types (e.g., pedestrian trajectory anomalies
rather than violent actions) and different video characteristics, remains to
be verified. The thesis is scoped to violence-oriented anomalies by design,
and its conclusions about skeleton complementarity are plausibly strongest
exactly there --- categories dominated by human motion --- so transfer to
appearance-dominated anomaly regimes should not be assumed.

Additionally, our corruption conditions for the robustness study are
synthetically generated following the ImageNet-C
protocol~\cite{hendrycks2019imagenetc}. Real-world distribution shifts
involving camera aging, environmental changes, compression pipelines, or
cross-city deployment may exhibit different characteristics, and the
transductive protocol discussed above compounds this caveat: the robustness
results should be read as a controlled study of the routing mechanism, not as
a field trial.

\section{Future Work}
\label{sec:lim-future}

The thesis infrastructure --- cached features, tracked run index, fixed
splits, and a sub-five-minute training loop --- leaves several follow-up
tracks ready to run. We list them as forward pointers only; none of the
numbers below are results of this thesis.

\textbf{Temporal modeling over snippets.} The fusion head scores each snippet
independently; the analysis in Chapter~\ref{ch:07} attributes the residual gap
to CLIP-TSA largely to its temporal self-attention machinery. A lightweight
temporal module (a TCN or single transformer layer over the snippet sequence)
is the largest single untapped headroom, but it changes the headline
architecture and would cascade through every table, so it was deliberately
deferred to a thesis extension rather than bolted on late.

\textbf{XD-Violence test-score smoothing (verified, deliberately not
adopted).} A simple post-hoc temporal smoothing of XD-Violence test scores
over a 64-frame window --- matching the native snippet length --- was verified
to add $+0.57$ points AP ($78.72 \to 79.29$, consistent across all three
seeds).\footnote{Tracked in the repository review artifact
\texttt{.planning/REVIEW-2026-06-10.md}, Section~5 (Pri-3), together with the
window-plateau analysis.} It is \emph{not} incorporated into any number in
this thesis: it was identified after the result freeze, and adopting a
post-hoc protocol change that moves the headline would undercut the
reproducibility discipline the thesis argues for. The follow-up is to
pre-register the smoothing window (showing the flat plateau around $w=64$
preempts the tuning objection) and adopt it in a future revision of the
pipeline.

\textbf{Skeleton-head ensemble routing.} The current reweighting routes
\emph{within} the frozen gated-fusion head. Routing instead to a separately
trained skeleton-only head (a reliability-routed score-level ensemble) is
estimated to add another 1--2 points to the corrupted-AUC mean, because the
skeleton stream's clean cache remains intact under the corruption families
where the fused head bottoms out; it would also likely convert the single-seed
SigLIP2-Base over-routing disclosure into a positive result.

\textbf{Streaming reliability estimation.} The routed signal is two scalars
per condition, which should converge quickly over a stream of test snippets. A
causal, streaming estimate of the routing weight would discharge the
transductive limitation above and upgrade the claim class from transductive to
online; until that experiment is run, no online claim is made.

\textbf{JPEG-family rescue diagnostic.} The reweighting is exactly zero on the
JPEG family by gate construction, because both streams degrade there --- the
largest untouched pool in the corruption study. Corrupted skeleton caches for
this family already exist, so a skeleton-only diagnostic would settle, with a
pre-registered kill criterion, whether any routing-based mechanism can help
where the skeleton stream itself is damaged.

\textbf{Text-prompt alignment branch.} The vision-language text tower is
entirely unused in this thesis. A lightweight VadCLIP-style alignment branch
over the cached joint-space features is the natural way to test whether the
text modality's gains survive the frozen-backbone constraint, and is the
flagship candidate for a follow-on extension.

\textbf{Adapter-based test-time adaptation.} Given the LayerNorm barrier
identified in Chapter~\ref{ch:07}, adaptation mechanisms that alter the
features themselves --- LoRA-style low-rank adapters or visual prompt tuning
--- are the principled next step for test-time robustness, as they sidestep
the affine-parameter bottleneck entirely.

\textbf{Audio for XD-Violence.} XD-Violence provides audio annotations that
our visual-skeleton pipeline does not exploit; audio cues are highly
discriminative for categories such as Shooting and Explosion, and the fusion
architecture extends naturally to a third stream.

\textbf{Additional backbones and datasets.} The frozen design makes new
vision-language backbones cheap to evaluate as they are released, and
extending the study to ShanghaiTech or CUHK Avenue would test how much of the
skeleton-complementarity finding is specific to violence-oriented anomalies.
